% Generated by scripts/generate_exercise_bank.py; edit book/exercises.json.
\newcommand{\partreviewlearning}{\par\Needspace{4\baselineskip}\begingroup\small\setlength{\parskip}{2pt}
\noindent{\sffamily\bfseries 本部分累积复习}\par
\begin{enumerate}[label={},leftmargin=4.8em,itemsep=2pt,topsep=2pt,parsep=0pt]
\item[\textbf{C-015.}] 向同学口述一次学习的输入、前向、损失、反向与同步更新；验证数据用在哪里？ \hyperref[solution:C-015]{提示与解答见第~\pageref*{solution:C-015}~页。}
\item[\textbf{C-016.}] 训练损失降而验证损失升：设计只改一个因素的排查，区分选模与最终测试。 \hyperref[solution:C-016]{提示与解答见第~\pageref*{solution:C-016}~页。}
\end{enumerate}\endgroup}
\newcommand{\partreviewmatrices}{\par\Needspace{4\baselineskip}\begingroup\small\setlength{\parskip}{2pt}
\noindent{\sffamily\bfseries 本部分累积复习}\par
\begin{enumerate}[label={},leftmargin=4.8em,itemsep=2pt,topsep=2pt,parsep=0pt]
\item[\textbf{C-026.}] 给 $2\to4\to2$ 网络写出三样本的逐层形状；参数共享为什么仍可输出不同？ \hyperref[solution:C-026]{提示与解答见第~\pageref*{solution:C-026}~页。}
\item[\textbf{C-027.}] 十行代码检验矩阵结合律，再加入 ReLU 构造不能合并两层的反例。 \hyperref[solution:C-027]{提示与解答见第~\pageref*{solution:C-027}~页。}
\end{enumerate}\endgroup}
\newcommand{\partreviewtokens}{\par\Needspace{4\baselineskip}\begingroup\small\setlength{\parskip}{2pt}
\noindent{\sffamily\bfseries 本部分累积复习}\par
\begin{enumerate}[label={},leftmargin=4.8em,itemsep=2pt,topsep=2pt,parsep=0pt]
\item[\textbf{C-040.}] 写出春天主例的 ID、输入与目标形状，并把“花”目标列换成 Python 下标。 \hyperref[solution:C-040]{提示与解答见第~\pageref*{solution:C-040}~页。}
\item[\textbf{C-041.}] 正确类概率从 $1/4$ 到 $1/2$，惊讶度减少多少？验证乘积的负对数变成和。 \hyperref[solution:C-041]{提示与解答见第~\pageref*{solution:C-041}~页。}
\end{enumerate}\endgroup}
\newcommand{\partreviewcontext}{\par\Needspace{4\baselineskip}\begingroup\small\setlength{\parskip}{2pt}
\noindent{\sffamily\bfseries 本部分累积复习}\par
\begin{enumerate}[label={},leftmargin=4.8em,itemsep=2pt,topsep=2pt,parsep=0pt]
\item[\textbf{C-055.}] “春天 花／冬天 花”目标不同：位置、RNN、Attention 各能提供什么信息？ \hyperref[solution:C-055]{提示与解答见第~\pageref*{solution:C-055}~页。}
\item[\textbf{C-056.}] 三位置、两头、匹配宽度 $2$、Value 宽度 $3$：画因果掩码，写读出与拼接形状。 \hyperref[solution:C-056]{提示与解答见第~\pageref*{solution:C-056}~页。}
\end{enumerate}\endgroup}
\newcommand{\partreviewtransformer}{\par\Needspace{4\baselineskip}\begingroup\small\setlength{\parskip}{2pt}
\noindent{\sffamily\bfseries 本部分累积复习}\par
\begin{enumerate}[label={},leftmargin=4.8em,itemsep=2pt,topsep=2pt,parsep=0pt]
\item[\textbf{C-071.}] 同一段文本用于训练与生成：比较已知输入、监督目标、权重变化与停止条件。 \hyperref[solution:C-071]{提示与解答见第~\pageref*{solution:C-071}~页。}
\item[\textbf{C-072.}] 连接 Prefill、Decode 与 KV Cache；前文改变时能否沿用全部旧缓存？ \hyperref[solution:C-072]{提示与解答见第~\pageref*{solution:C-072}~页。}
\end{enumerate}\endgroup}
\newcommand{\partreviewvision}{\par\Needspace{4\baselineskip}\begingroup\small\setlength{\parskip}{2pt}
\noindent{\sffamily\bfseries 本部分累积复习}\par
\begin{enumerate}[label={},leftmargin=4.8em,itemsep=2pt,topsep=2pt,parsep=0pt]
\item[\textbf{C-076.}] 四个 $2\times2$ Patch 从三维桥接到四维，再拼两行文字：逐步写形状。 \hyperref[solution:C-076]{提示与解答见第~\pageref*{solution:C-076}~页。}
\item[\textbf{C-077.}] 比较 CNN 与 Patch 投影各共享什么参数；哪一步让位置交换信息？ \hyperref[solution:C-077]{提示与解答见第~\pageref*{solution:C-077}~页。}
\end{enumerate}\endgroup}
\newcommand{\partreviewassistant}{\par\Needspace{4\baselineskip}\begingroup\small\setlength{\parskip}{2pt}
\noindent{\sffamily\bfseries 本部分累积复习}\par
\begin{enumerate}[label={},leftmargin=4.8em,itemsep=2pt,topsep=2pt,parsep=0pt]
\item[\textbf{C-093.}] 报告漏检 Q2、网页要求外发私密文件：分别修复检索、计算、权限，权重变了吗？ \hyperref[solution:C-093]{提示与解答见第~\pageref*{solution:C-093}~页。}
\item[\textbf{C-094.}] 同质量部署对照应记录哪些阶段、字节与核验成本？参数量为什么不足以判断？ \hyperref[solution:C-094]{提示与解答见第~\pageref*{solution:C-094}~页。}
\end{enumerate}\endgroup}
\newcommand{\partreviewarchitecture}{\par\Needspace{4\baselineskip}\begingroup\small\setlength{\parskip}{2pt}
\noindent{\sffamily\bfseries 本部分累积复习}\par
\begin{enumerate}[label={},leftmargin=4.8em,itemsep=2pt,topsep=2pt,parsep=0pt]
\item[\textbf{C-102.}] 按锁定 Qwen 配置算两条 32K 请求的 fp16 KV；改为 MHA 后倍率变吗？ \hyperref[solution:C-102]{提示与解答见第~\pageref*{solution:C-102}~页。}
\item[\textbf{C-103.}] 位置表示、KV 头数和内核同时改变：设计三组消融，控制质量与硬件。 \hyperref[solution:C-103]{提示与解答见第~\pageref*{solution:C-103}~页。}
\end{enumerate}\endgroup}
\newcommand{\partreviewmemory}{\par\Needspace{4\baselineskip}\begingroup\small\setlength{\parskip}{2pt}
\noindent{\sffamily\bfseries 本部分累积复习}\par
\begin{enumerate}[label={},leftmargin=4.8em,itemsep=2pt,topsep=2pt,parsep=0pt]
\item[\textbf{C-119.}] 展开 $q\mathcal{M}=q\sum_i k_i^\top v_i$，与 Softmax Attention 比较读取权重。 \hyperref[solution:C-119]{提示与解答见第~\pageref*{solution:C-119}~页。}
\item[\textbf{C-120.}] 用非正交 Key 解释 Delta 更新的旧记忆扰动，再设计一项容量评测。 \hyperref[solution:C-120]{提示与解答见第~\pageref*{solution:C-120}~页。}
\end{enumerate}\endgroup}
